\documentclass[10pt,letterpaper]{article}
\usepackage[T1]{fontenc}
\usepackage[utf8]{inputenc}
\usepackage[sfdefault]{sourcesanspro}
\usepackage{inconsolata}
\usepackage[margin=0.84in,headheight=15pt]{geometry}
\usepackage{microtype}
\usepackage{xcolor}
\usepackage{booktabs,longtable,tabularx,array}
\usepackage{enumitem}
\usepackage{fancyhdr,lastpage}
\usepackage[most]{tcolorbox}
\usepackage{listings}
\usepackage{titlesec}
\usepackage{hyperref}
\usepackage{amsmath,amssymb}
\usepackage{multicol}
\definecolor{navy}{HTML}{15324C}
\definecolor{teal}{HTML}{087E83}
\definecolor{slate}{HTML}{536879}
\definecolor{paper}{HTML}{F2F6F8}
\definecolor{warm}{HTML}{E8F4F3}
\hypersetup{colorlinks=true,urlcolor=teal,linkcolor=navy,pdfauthor={Software Engineering Reference Suite},pdfsubject={MISRA C++:2023 engineering guidance}}
\titleformat{\section}{\Large\bfseries\color{navy}}{\thesection}{0.7em}{}
\titleformat{\subsection}{\large\bfseries\color{teal}}{\thesubsection}{0.6em}{}
\titleformat{\subsubsection}{\normalsize\bfseries\color{navy}}{\thesubsubsection}{0.5em}{}
\setlist[itemize]{leftmargin=*,topsep=2pt,itemsep=2pt}
\setlist[enumerate]{leftmargin=*,topsep=2pt,itemsep=3pt}
\setlength{\parskip}{5pt}
\setlength{\parindent}{0pt}
\setlength{\emergencystretch}{2em}
\pagestyle{fancy}
\fancyhf{}
\fancyhead[L]{\small\color{slate}MISRA C++:2023 Engineering Suite}
\fancyhead[R]{\small\color{slate}\nouppercase{\leftmark}}
\fancyfoot[L]{\scriptsize\color{slate}Independent engineering guidance; see licensed source}
\fancyfoot[R]{\scriptsize\color{slate}\thepage\ / \pageref{LastPage}}
\newtcolorbox{sourcebox}{colback=paper,colframe=navy!30,boxrule=0.45pt,arc=1.6pt,breakable}
\newtcolorbox{practicebox}{colback=warm,colframe=teal!40,boxrule=0.45pt,arc=1.6pt,breakable}
\newcommand{\SourceNote}[1]{\begin{sourcebox}\textbf{SOURCE BASIS / NORMATIVE POINTER}\quad #1\end{sourcebox}}
\newcommand{\PracticeNote}[1]{\begin{practicebox}\textbf{ENGINEERING RECOMMENDATION / NOT MISRA TEXT}\quad #1\end{practicebox}}
\newcommand{\Record}[1]{\textsf{\bfseries #1}}
\lstdefinestyle{cppsuite}{language=C++,basicstyle=\ttfamily\footnotesize,keywordstyle=\color{navy}\bfseries,commentstyle=\color{slate},stringstyle=\color{teal},breaklines=true,breakatwhitespace=true,columns=fullflexible,frame=single,framerule=0.3pt,rulecolor=\color{navy!35},backgroundcolor=\color{paper},showstringspaces=false,tabsize=2,xleftmargin=0.5em,xrightmargin=0.5em,aboveskip=4pt,belowskip=4pt}
\lstset{style=cppsuite}
\newcommand{\SuiteTitle}[3]{%
\begin{titlepage}\thispagestyle{empty}
\vspace*{1.5cm}
{\color{teal}\rule{2.4cm}{2pt}}\par\vspace{1cm}
{\Huge\bfseries\color{navy}MISRA C++:2023\\[0.4em]#1}\par\vspace{0.7cm}
{\Large\color{slate}#2}\par\vspace{1.2cm}
\begin{sourcebox}\textbf{Source baseline:} MISRA C++:2023, Guidelines for the use of C++17 in critical systems, October 2023 (user-supplied licensed copy). This is an independent companion, not an official MISRA publication, rule replacement, certified checker, or claim of project compliance.\end{sourcebox}
\vfill{\large\color{navy}Technical edition \textbullet\ 16 October 2026}\par\smallskip
{\small\color{slate}#3}\par\end{titlepage}
\tableofcontents\clearpage
}
\newcommand{\StandardsBoundary}{%
\section*{How to read normative status}\addcontentsline{toc}{section}{How to read normative status}
\textbf{Normative source:} the licensed MISRA C++:2023 guideline text, including any applicable \emph{Amplification} and \emph{Exception}. A guideline \emph{headline} alone is not a complete compliance test. The source's Section~3 defines classifications, categories, decidability, analysis scope, and applicability.\par
\textbf{Independent guidance:} all review questions, programming patterns, lifecycle processes, sample thresholds, ownership assignments, controls, and CI configurations in this suite are \emph{engineering recommendations}, not statements issued by MISRA.\par
\textbf{Critical limitation:} no individual example is asserted compliant with the whole standard. For formal compliance claims, follow the applicable edition of \emph{MISRA Compliance} as the source itself instructs, with a documented compliance plan and supporting evidence.\par
}

\begin{document}
\SuiteTitle{Hands-On Implementation Handbook}{Eight prerequisite-linked engineering labs and capstone acceptance}{Volume 10 / Engineering Labs}
\StandardsBoundary

\section{Lab environment and conventions}
Required prerequisites: C++17 compiler (GCC/Clang or target compiler), build system, unit-test runner, source-controlled repository, authorized access to the licensed MISRA source and a project-qualified rule checker or a documented manual review alternative. The fixture folder provides 22 contrasted examples.

Do not execute arbitrary noncompliant specimen code; use syntax-only compilation, static analysis or isolated test harnesses. Noncompliant examples include undefined behavior risks. Expected completion artifacts are suggested audit evidence, not MISRA-defined mandatory form names.


\section{Lab 01: Compile database and baseline}
\textbf{Prerequisites:} C++17 compiler, CMake.\\\textbf{Objective:} Capture toolchain and variant facts before any analysis.

\subsection{Procedure}
\begin{enumerate}
\item Create a minimal two-file C++17 project and generate compile\_commands.json.
\item Record preprocessor macro defines, ABI, standard-library implementation and compiler version.
\item Compile both debug and release variants; compare translation-unit coverage.
\item Document the exact software item and product configurations under review.
\end{enumerate}

\subsection{Acceptance evidence}
A reproducible configuration manifest, source revision and two build logs; no untracked translation units.

\textbf{Discussion question:} Which part of the result is demonstrated by your artifacts, and which requires reading normative amplification, exception or external compliance guidance?\par

\section{Lab 02: Undefined behavior and lifetime}
\textbf{Prerequisites:} Lab 01; pointer and reference fundamentals.\\\textbf{Objective:} Analyze Rule 6.8.2 and Rule 11.6.2 against real control flow.

\subsection{Procedure}
\begin{enumerate}
\item Inspect examples/noncompliant/6\_8\_2\_lifetime.cpp and 11\_6\_2\_init.cpp.
\item Identify lifetime end or uninitialized read on the violating path.
\item Apply the alternative designs, run a C++17 syntax check and source-guideline scan.
\item Review source-defined amplification/exception for both guideline IDs.
\end{enumerate}

\subsection{Acceptance evidence}
A before/after code diff, two rule citations, path argument and checker observations.

\textbf{Discussion question:} Which part of the result is demonstrated by your artifacts, and which requires reading normative amplification, exception or external compliance guidance?\par

\section{Lab 03: Integral conversions and null-pointer constants}
\textbf{Prerequisites:} Lab 02; integral promotions.\\\textbf{Objective:} Develop a review argument for Rules 7.0.1, 7.0.2 and 7.11.1.

\subsection{Procedure}
\begin{enumerate}
\item Construct a table of static source and destination types.
\item Review the conversion and nullptr fixture pairs.
\item Add boundary and signedness-focused unit tests.
\item Explain which findings cannot be judged solely from a compiler warning.
\end{enumerate}

\subsection{Acceptance evidence}
Type-transition matrix, correction patch and targeted tests.

\textbf{Discussion question:} Which part of the result is demonstrated by your artifacts, and which requires reading normative amplification, exception or external compliance guidance?\par

\section{Lab 04: Flow graphs, unreachable code and recursion}
\textbf{Prerequisites:} Lab 03; CFG and call graphs.\\\textbf{Objective:} Distinguish syntactic unreachable blocks from path-infeasible branches.

\subsection{Procedure}
\begin{enumerate}
\item Draw the CFG of Rule 0.0.1 fixture and mark return edges.
\item Construct one example containing constexpr if and explain source applicability.
\item Build a call graph for Rule 8.2.10 recursive fixture.
\item Rewrite the recursive algorithm iteratively and test equivalent behavior for a safe input range.
\end{enumerate}

\subsection{Acceptance evidence}
Annotated CFG, call graph, source-exception notes and executable equivalence tests.

\textbf{Discussion question:} Which part of the result is demonstrated by your artifacts, and which requires reading normative amplification, exception or external compliance guidance?\par

\section{Lab 05: Class and ownership patterns}
\textbf{Prerequisites:} Lab 04; RAII, constructors.\\\textbf{Objective:} Apply Rules 15.1.3, 21.6.2, 23.11.1 and 28.6.3.

\subsection{Procedure}
\begin{enumerate}
\item Model a small resource-managing class and record ownership invariants.
\item Make single-argument conversion construction explicit.
\item Replace unmanaged allocations with type-safe ownership.
\item Exercise move construction and establish a post-move contract.
\end{enumerate}

\subsection{Acceptance evidence}
Class contract, lifecycle test log and rationale for each changed expression.

\textbf{Discussion question:} Which part of the result is demonstrated by your artifacts, and which requires reading normative amplification, exception or external compliance guidance?\par

\section{Lab 06: Headers, macro restrictions and build variants}
\textbf{Prerequisites:} Lab 01; preprocessors.\\\textbf{Objective:} Demonstrate Rules 19.0.2 and 19.2.1 plus variant analysis.

\subsection{Procedure}
\begin{enumerate}
\item Remove function-like macros using constexpr functions where practical.
\item Add include guards to a header used from two independent translation units.
\item Compare preprocessor output for two macro configurations.
\item Record whether system-scope evidence covers both shipping variants.
\end{enumerate}

\subsection{Acceptance evidence}
Preprocessed snapshots, two clean builds and feature-configuration matrix.

\textbf{Discussion question:} Which part of the result is demonstrated by your artifacts, and which requires reading normative amplification, exception or external compliance guidance?\par

\section{Lab 07: Analyzer evaluation and disposition}
\textbf{Prerequisites:} Labs 01--06.\\\textbf{Objective:} Develop a check-specific tool qualification mini-suite.

\subsection{Procedure}
\begin{enumerate}
\item Run each pair in the fixture suite through the selected MISRA analyzer.
\item Map vendor checker IDs to the exact 179 source guideline IDs.
\item Write source-exception and false-positive demonstrations for one rule.
\item Complete one hypothetical Required deviation draft; do not approve it without authority.
\end{enumerate}

\subsection{Acceptance evidence}
Coverage matrix, false-positive analysis and draft deviation record.

\textbf{Discussion question:} Which part of the result is demonstrated by your artifacts, and which requires reading normative amplification, exception or external compliance guidance?\par

\section{Lab 08: Release-assurance capstone}
\textbf{Prerequisites:} Labs 01--07.\\\textbf{Objective:} Produce a bounded compliance evidence package.

\subsection{Procedure}
\begin{enumerate}
\item Define an exact simulated shipping configuration.
\item Collect compile database, unit tests, analysis output and rule coverage.
\item Adjudicate all findings by category and document unresolved uncertainty.
\item Run the example record-consistency policy gate with a complete mock coverage map.
\item Prepare a clearly labeled simulated compliance assessment for independent review.
\end{enumerate}

\subsection{Acceptance evidence}
Traceable evidence bundle with explicit qualifications; no unsupported compliance claim.

\textbf{Discussion question:} Which part of the result is demonstrated by your artifacts, and which requires reading normative amplification, exception or external compliance guidance?\par

\section{Capstone grading matrix}
\small\setlength{\tabcolsep}{3pt}\renewcommand{\arraystretch}{1.15}
\begin{longtable}{p{0.26\linewidth}p{0.12\linewidth}p{0.5\linewidth}}
\toprule
Criterion & Weight & Pass condition\\ \midrule
\endfirsthead
\toprule
Criterion & Weight & Pass condition\\ \midrule
\endhead
Source interpretation & 25\% & Exact guideline citations and correct category/exception differentiation.\\
C++ correctness & 20\% & Corrected code builds and behavior is tested.\\
Analysis integrity & 20\% & Compiler/context and tool limitations visible.\\
Disposition & 20\% & Mandatory/Required/Advisory workflow distinctions applied.\\
Evidence and independence & 15\% & Independent review questions and artifact traceability included.\\
\bottomrule
\end{longtable}
\normalsize
Suggested educational passing score: 85%; require 100% correct answers on the Mandatory category/deviation rule. These values are local curriculum recommendations and should not be quoted as MISRA requirements.


\end{document}
